\documentclass[11pt,a4paper]{article}

\usepackage[T1]{fontenc}
\usepackage[utf8]{inputenc}
\usepackage[spanish,es-noshorthands]{babel}
\usepackage{lmodern}
\usepackage{amsmath,amssymb,amsthm}
\usepackage[margin=2.5cm]{geometry}
\usepackage{booktabs}
\usepackage{microtype}
\usepackage{enumitem}
\usepackage[colorlinks=true,linkcolor=black,citecolor=black,urlcolor=blue]{hyperref}

\newcommand{\pc}{\ifmmode\text{\,\%}\else\,\%\fi}
\newcommand{\MER}{\mathrm{MER}}
\newcommand{\dd}{\,\mathrm{d}}
\newcommand{\p}{\partial}
\newcommand{\Rv}{R_v}

\title{\vspace{-1.4cm}\textbf{Formulación matemática del modelo bifásico\\ de ascenso de magma}\\[0.25em]
\large El operador directo de \texttt{RIconduitex5\_5}}
\author{Ignacio Gómez Gómez}
\date{\today}

\begin{document}
\maketitle
\vspace{-0.9cm}

\begin{abstract}\noindent
Se escribe el modelo de conducto que ya está implementado
(\texttt{RIconduitex5\_5.py}), no un recorte. Es un flujo bifásico
estacionario unidimensional, isotermo, de sección constante, con dos
velocidades y arrastre, en el marco de Slezin (2003) y Kozono \&
Koyaguchi (2009). El estado es
$\mathbf{y}=(P,\phi,N,\xi,u_m,u_g)$. Las dos velocidades son
algebraicas; $P$, $\phi$, $N$ y $\xi$ cumplen un DAE. El caudal másico
$q$ se determina por un problema de tiro. Ése es el operador directo
que se quiere invertir.
\end{abstract}

\tableofcontents
\vspace{0.6em}

\section{Dominio, fases e hipótesis}

El conducto es el intervalo $z\in[H,0]$, con $z$ creciente hacia la
superficie ($H<0$ es el techo de la cámara). La sección es constante:
cilindro de radio $R$, o dique de semi-ancho $R$. El flujo es
estacionario e isotermo. Hay dos fases con \textbf{velocidad propia}:

\begin{center}
\begin{tabular}{@{}lll@{}}
\toprule
fase & densidad & velocidad \\
\midrule
fundido (silicato $+$ cristales $+$ agua disuelta) & $\rho_m$ (constante) & $u_m(z)$ \\
gas (vapor de agua) & $\rho_g=P/(\Rv T)$ & $u_g(z)$ \\
\bottomrule
\end{tabular}
\end{center}

$\phi(z)$ es la fracción volumétrica de gas. La fracción volumétrica
de fundido es $1-\phi$. Los cristales viven en el fundido: $\xi(z)$ es
su fracción en la fase condensada. $N(z)$ es el número de burbujas por
unidad de volumen de fundido.

Hipótesis que el código impone y que hay que tener a la vista:

\begin{itemize}[leftmargin=1.2em,itemsep=0.2em]
\item una sola presión (las dos fases están en equilibrio mecánico
  local: no hay $P_g\neq P_m$);
\item $\rho_m$ se calcula una vez, a la presión y al agua disuelta de
  la base, y se mantiene constante (Bottinga \& Weill 1970); el
  comentario de \texttt{calbuco2015d.py} deja explícito que $\xi$
  \emph{aún no} entra en esa densidad;
\item el agua disuelta está en equilibrio con $P$ (no hay ecuación de
  retraso para $c_d$);
\item $T$ es uniforme;
\item el área transversal no depende de $z$.
\end{itemize}

Caso base (Calbuco 2015, \texttt{calbuco2015d.py}):
$H=-7000\,\mathrm{m}$, $R=16\,\mathrm{m}$,
$\Delta P=5\,\mathrm{MPa}$, $c_0=0.04$, $T=1243.15\,\mathrm{K}$,
$\xi_0=0.25$.

\section{Variables de estado}

El estado en cada $z$ es el vector de seis componentes
\begin{equation}
\label{eq:estado}
\mathbf{y}(z)=\bigl(P(z),\;\phi(z),\;N(z),\;\xi(z),\;u_m(z),\;u_g(z)\bigr).
\end{equation}
Las dos últimas son \emph{algebraicas}: quedan determinadas por la
conservación de masa una vez conocidos $(P,\phi,\xi)$ y el caudal
$q$. Las cuatro primeras son diferenciales. El sistema es un DAE de
índice~1,
\begin{equation}
\label{eq:DAE}
M\,\mathbf{y}'=F(\mathbf{y};q),\qquad
M=\mathrm{diag}(1,1,1,1,0,0).
\end{equation}

\section{Conservación de masa}

Como el área es constante y el flujo es estacionario, el caudal másico
por unidad de área es un escalar,
\begin{equation}
\label{eq:q}
q
=\rho_m(1-\phi)\,u_m
+\rho_g\,\phi\,u_g
=\mathrm{const}.
\end{equation}
Ese $q$ es el $\MER$ por unidad de área; para un cilindro,
$\MER=\pi R^2 q$. Se descompone en la fracción másica de gas exsuelto
\begin{equation}
\label{eq:n}
n
:=\frac{\rho_g\,\phi\,u_g}{q},
\qquad
1-n
=\frac{\rho_m(1-\phi)\,u_m}{q},
\end{equation}
y se despejan las velocidades
\begin{equation}
\label{eq:vel}
u_m
=\frac{q\,(1-n)}{(1-\phi)\,\rho_m},
\qquad
u_g
=\frac{q\,n}{\phi\,\rho_g}.
\end{equation}
Esto es exactamente \texttt{result[4]} y \texttt{result[5]} del
código. $n$ no es una incógnita extra: es una función algebraica de
$(P,\xi)$, dada por la solubilidad (sección~\ref{sec:sol}).

En la base, si aún no hay gas exsuelto, $n=0$, $\phi=0$ y
$u_m=u_g=v_{\mathrm{in}}$, con $q=v_{\mathrm{in}}\rho_m$. Si ya hay
gas,
\begin{equation}
\rho_{t,i}
=\frac{P_i\rho_m}{\rho_m \Rv T\,n_i+P_i(1-n_i)},
\qquad
\phi_i
=\Biggl(1+\frac{P_i}{n_i\Rv T}\frac{1-n_i}{\rho_m}\Biggr)^{-1},
\qquad
q=v_{\mathrm{in}}\rho_{t,i}.
\end{equation}

\section{Conservación de momentum}
\label{sec:mom}

Con $z$ creciente hacia arriba, el balance de momentum de cada fase
es el de Kozono \& Koyaguchi (2009):
\begin{align}
\rho_m(1-\phi)\,u_m\frac{\dd u_m}{\dd z}
&=-(1-\phi)\frac{\dd P}{\dd z}
-\rho_m(1-\phi)\,g
+F_{mg}-F_{mw},
\label{eq:mom-m}\\[0.3em]
\rho_g\phi\,u_g\frac{\dd u_g}{\dd z}
&=-\phi\frac{\dd P}{\dd z}
-\rho_g\phi\,g
-F_{mg}-F_{gw}.
\label{eq:mom-g}
\end{align}
$F_{mw}$ es la fricción pared--fundido, $F_{gw}$ la pared--gas y
$F_{mg}$ el arrastre gas--fundido (fuerza por unidad de volumen
\emph{sobre el fundido}; sobre el gas va con signo contrario).

Las velocidades~\eqref{eq:vel} dependen de $P$ y $\phi$ (y de $\xi$,
que entra en $n$). Al sustituir $\dd u_m/\dd z$ y $\dd u_g/\dd z$ en
\eqref{eq:mom-m}--\eqref{eq:mom-g} queda un sistema lineal $2\times2$
para $(\dd P/\dd z,\,\dd\phi/\dd z)$. El código lo resuelve en la
forma
\begin{equation}
\label{eq:22}
\frac{\dd\phi}{\dd z}
=\frac{-e\,b+d\,f}{ad-bc},
\qquad
\frac{\dd P}{\dd z}
=\frac{-e\,a+f\,c}{ad-bc},
\end{equation}
con los coeficientes de \texttt{momenteq}
\begin{align}
a&=\rho_g u_g^2, &
b&=\phi-\frac{u_g^2\phi}{\Rv T}+\frac{\p n}{\p P}\,q\,u_g,
\label{eq:coef}\\
c&=\rho_m u_m^2, &
d&=\frac{\p n}{\p P}\,q\,u_m-(1-\phi),\nonumber\\
e&=-\rho_m(1-\phi)\,g+F_{mg}-F_{mw}, &
f&=\rho_g\phi\,g+F_{mg}+F_{gw}.\nonumber
\end{align}
La cantidad $f$ es menos el lado derecho de~\eqref{eq:mom-g} sin el
término de presión: el signo de $F_{mg}$ en $e$ y en $f$ es el mismo
\emph{después} de haber pasado~\eqref{eq:mom-g} al otro lado. El
arrastre, en la física, es opuesto en las dos fases.

Los términos con $\p n/\p P$ vienen de derivar~\eqref{eq:vel} respecto
de $P$ a través de $n=n(P,\xi)$ y de $\rho_g=P/(\Rv T)$. Sin ellos el
sistema no sería consistente con la exsolución.

\section{Solubilidad y gas exsuelto}
\label{sec:sol}

Ley de Henry,
\begin{equation}
\label{eq:henry}
c_s(P)=C_1 P^\beta,
\qquad
C_1=4.11\times10^{-6},
\quad
\beta=\tfrac12,
\end{equation}
con $P$ en pascales y $c_s$ en fracción másica del fundido. El agua
total del sistema es $c_0$. Los cristales no disuelven agua: el
contenido de agua ``disponible'' en la fase condensada se corrige por
$\xi$. El código usa
\begin{equation}
\label{eq:test}
\mathrm{test}
=(1-\xi_0)\,c_0-(1-\xi)\,C_1 P^\beta.
\end{equation}
Si $\mathrm{test}\le0$ el magma está subsaturado: $n=0$. Si
$\mathrm{test}>0$,
\begin{equation}
\label{eq:fg}
n(P,\xi)
=\frac{(1-\xi_0)\,c_0-(1-\xi)\,C_1 P^\beta}{1-C_1 P^\beta}.
\end{equation}
La derivada $\p n/\p P$ que entra en~\eqref{eq:coef} es la de la
regla del cociente sobre~\eqref{eq:fg}, e incluye
$\p\xi/\p P$ vía la cinética (sección~\ref{sec:xtal}). Con
$A=(1-\xi_0)c_0-(1-\xi)C_1P^\beta$ y $B=1-C_1P^\beta$,
\begin{equation}
\frac{\p n}{\p P}
=\frac{A_P B-A B_P}{B^2},
\qquad
A_P=\frac{\p\xi}{\p P}\,C_1P^\beta-(1-\xi)C_1\beta P^{\beta-1},
\qquad
B_P=-C_1\beta P^{\beta-1}.
\end{equation}
En el código es \texttt{dfgdp}.

\section{Cristales}
\label{sec:xtal}

Hay una fracción inicial de fenocristales $\xi_0$ y un techo
$\xi_{\max}$. La fracción de equilibrio a la presión $P$ se interpola
con el agua ya exsuelta,
\begin{equation}
\label{eq:xteo}
f_2=\max\Biggl(0,\;
\frac{(1-\xi_0)c_0-(1-\xi)C_1 P^\beta}
{(1-\xi_0)c_0-(1-\xi_{\max})C_1 P_{\mathrm{atm}}^\beta}\Biggr),
\qquad
\xi_{\mathrm{eq}}
=\xi_0+(\xi_{\max}-\xi_0)f_2,
\end{equation}
y se cierra con $f_3=\max(0,\,1-\xi/\xi_{\mathrm{eq}})$. La tasa es
de primer orden, con tiempo característico $\tau_c$
(\texttt{tcar}${}=2\,\mathrm{h}$):
\begin{equation}
\label{eq:dxdt}
\frac{\dd\xi}{\dd t}
=\max\Biggl(0,\;
\frac{(\xi_{\max}-\xi_0)\,f_2 f_3}{\tau_c}\Biggr).
\end{equation}
A lo largo del conducto, $\dd\xi/\dd z=(\dd\xi/\dd t)/u_m$. El código
usa esa forma en el régimen $n_{\mathrm{eq}}=3$; en los regímenes
$1$ y $2$ escribe $\dd\xi/\dd z$ \emph{sin} dividir por $u_m$ (hay un
comentario ``Preguntar'' en \texttt{momenteq}). La forma
físicamente consistente es la primera.

\section{Burbujas: radio y coalescencia}

El radio de burbuja compatible con $\phi$ y $N$ es
\begin{equation}
\label{eq:rb}
r_b
=\Biggl(\frac{\phi}{\tfrac43\pi N(1-\phi)}\Biggr)^{1/3},
\end{equation}
es decir, $N$ está medida por unidad de volumen de fundido:
$\phi=N\cdot(\tfrac43\pi r_b^3)\cdot(1-\phi)$.

La coalescencia sigue una ley de tipo Slezin. Mientras
$r_b<\tfrac12 R$ y $\phi<\phi_{\mathrm{crit}}$,
\begin{equation}
\label{eq:dNdt}
\begin{aligned}
\frac{\dd N}{\dd t}
&=-N^{2/3}\,(1-\phi)^{-1/3}
\cdot\frac19\frac{(\rho_m-\rho_g)g}{\mu}
\cdot\Biggl(\frac{3\phi}{4\pi}\Biggr)^{2/3}\\
&\qquad\cdot(F_1^2-F_2^2)
\cdot\Biggl(1-(F_1+F_2)
\Biggl(\frac{3\phi\cdot\pi/(6\phi_{\mathrm{crit}})}{4\pi}\Biggr)^{1/3}\Biggr)^{-1}\\
&\qquad\cdot F_c\,(1-\phi/\phi_{\mathrm{crit}})\,(R-r_b)/R,
\end{aligned}
\end{equation}
con $F_1=1.8$, $F_2=0.2$ (tamaños relativos de la población bimodal)
y $F_c\in[0,1]$ la eficiencia. Si no se cumple la guarda,
$\dd N/\dd t=0$. A lo largo del conducto,
\begin{equation}
\frac{\dd N}{\dd z}=\frac{1}{u_m}\frac{\dd N}{\dd t}.
\end{equation}
$N$ en la base se fija por la tasa de descompresión (Toramaru):
$\log_{10}N=1.5\log_{10}(\dot P)+5$, con
$\dot P=-(\dd P/\dd z)\,v_{\mathrm{in}}$, y un suelo $N=10^8$ si hay
exsolución ya en la entrada.

\section{Fuerzas}

\subsection{Fricción con la pared}

En el tramo viscoso $F_{gw}=0$ y
\begin{equation}
\label{eq:Fmw}
F_{mw}
=c_g\,\frac{\mu\,u_m}{R^2},
\qquad
c_g
=\begin{cases}
8 & \text{cilindro (Poiseuille)},\\
3 & \text{dique}.
\end{cases}
\end{equation}
Tras la fragmentación $F_{mw}=0$ y aparece
\begin{equation}
F_{gw}=0.01\,\frac{\rho_g\,|u_g|\,u_g}{4R}.
\end{equation}

\subsection{Arrastre gas--fundido}

Depende del régimen, marcado por la fracción de gas. Se definen dos
umbrales de permeabilidad $\phi_1<\phi_2$ y el umbral de
fragmentación $\phi_{\mathrm{crit}}$ (sección~\ref{sec:reg}).
$Re_b=2 r_b\rho_g|u_g-u_m|/\mu_g$ con $\mu_g=10^{-5}\,\mathrm{Pa\,s}$.

\paragraph{Régimen 1 ($\phi<\phi_1$).}
Stokes,
\begin{equation}
F_{mg}=\frac{3\mu}{r_b^2}\,(u_g-u_m)\,\phi(1-\phi).
\end{equation}

\paragraph{Régimen 2 ($\phi_1\le\phi<\phi_2$).}
Interpolación $\theta=(\phi-\phi_1)/(\phi_2-\phi_1)$. Si $Re_b>2200$,
entre Stokes y arrastre inercial; si no, entre Stokes y Darcy, con
permeabilidad
$k=0.131\,r_b^2(\phi-\phi_1+0.05)^{2.1}$:
\begin{equation}
F_{mg}
=\Biggl(\frac{\mu_g}{k}\Biggr)^\theta
\Biggl(\frac{3\mu}{r_b^2}\Biggr)^{1-\theta}
(u_g-u_m)\,\phi(1-\phi)
\end{equation}
(y el análogo inercial con $0.33\rho_g|u_g-u_m|/(4r_b)$ en lugar de
$\mu_g/k$).

\paragraph{Régimen 3 ($\phi_2\le\phi<\phi_{\mathrm{crit}}$).}
El extremo $\theta=1$ de lo anterior: Darcy o inercial, sin
mezclar con Stokes.

\paragraph{Régimen 4 ($\phi\ge\phi_{\mathrm{crit}}$, fragmentado).}
El fundido ya no es continuo. El arrastre es el de una partícula en
gas ($r_a=10^{-3}\,\mathrm{m}$, $c_d=0.8$), interpolado en una capa
de ancho $0.05$ en $\phi$:
\begin{equation}
F_{mg}
=\frac{3c_d}{8r_a}\,\rho_g\,|u_g-u_m|(u_g-u_m)\,\phi(1-\phi)
\end{equation}
cuando $\phi\ge\phi_{\mathrm{crit}}+0.05$.

\section{Reología}

La viscosidad del fundido es la de Giordano \emph{et al.} (2008),
función de la composición anhidra, del agua disuelta $c_d$ y de $T$.
Si el magma está saturado, $c_d=C_1 P^\beta$; si no, $c_d=c_0$. Los
cristales multiplican por un factor relativo
$\theta_\xi=\mathtt{fvrel}(\xi,\xi_0,\dot\gamma)$, con
$\dot\gamma=u_m/R$. El caso base usa el medio efectivo de dos sólidos
(\texttt{em2s}):
\begin{equation}
\theta_\xi
=\Bigl(1-\frac{\xi_0}{\phi_{\max,1}}\Bigr)^{-2.5}
\Bigl(1-\frac{\xi-\xi_0}{(1-\xi_0)\phi_{\max,2}}\Bigr)^{-2.5},
\end{equation}
con $\phi_{\max,i}=0.656\exp\bigl(-(\log_{10}r_i)^2/(2\cdot1.08^2)\bigr)$
y razones de aspecto $r_1=4$, $r_2=8$.

Las burbujas corrigen otra vez, a la manera de Llewellin, con el
número de capilaridad $Ca=r_b\mu_\ell\dot\gamma/3$:
\begin{equation}
\label{eq:mubub}
\mu
=\Biggl(\tfrac12(A-B)\bigl(1-\mathrm{erf}(c_1\log Ca+c_2)\bigr)+B\Biggr)
\theta_\xi\,\mu_{\mathrm{fundido}},
\end{equation}
donde $A=(1-\phi/\phi_\ast)^{-\phi_\ast}$,
$B=(1-\phi/\phi_\ast)^{5\phi_\ast/3}$, $\phi_\ast=\phi_{\mathrm{crit}}+0.05$,
$c_1=-0.2895\phi+0.8132$ y $c_2=\phi$.

\section{Regímenes, fragmentación y condiciones de borde}
\label{sec:reg}

El umbral de fragmentación no es un dato fijo. Se calcula un número
de capilaridad de referencia a $\phi\approx0.2$ y
\begin{equation}
\label{eq:phicrit}
\phi_{\mathrm{crit}}
=\frac{\ell_+-\ell_-}{2}\,\mathrm{erf}\bigl(\log_{10}Ca\bigr)
+\frac{\ell_++\ell_-}{2},
\end{equation}
con $\ell_+=0.785$, $\ell_-=0.525$. Los umbrales de permeabilidad
$\phi_1$, $\phi_2=\phi_1+0.01$ salen de la misma forma, con
$(\ell_+,\ell_-)=(0.40,0.15)$. El integrador marcha por tramos y
cambia de régimen al cruzar $\phi_1$, $\phi_2$ y
$\phi_{\mathrm{crit}}$. Tras la fragmentación se congelan $N$ y $\xi$,
y el estado diferencial se reduce a $(P,\phi)$.

En la base,
\begin{equation}
\label{eq:bc-base}
P(H)=P_{\mathrm{lit}}(H)+\Delta P
=\rho_{\mathrm{crust}}g|H|+\Delta P,
\qquad
\xi(H)=\xi_0,
\qquad
u_m(H)=u_g(H)=v_{\mathrm{in}}.
\end{equation}
$v_{\mathrm{in}}$ \textbf{no} es un dato: es el parámetro de tiro. En
la superficie se pide una de las dos condiciones de Yoshida \&
Koyaguchi (1999) / Kozono \& Koyaguchi (2009),
\begin{equation}
\label{eq:bc-top}
P(0)=P_{\mathrm{atm}}
\quad\text{o}\quad
u_g(0)\approx c_s:=\sqrt{\Rv T}
\end{equation}
(la segunda es el estrangulamiento; $c_s$ es la velocidad del sonido
isoterma del gas). El código acepta la solución si $z_{\mathrm{exit}}\ge-5\,\mathrm{m}$
y se cumple el choque o $P\approx P_{\mathrm{atm}}$ con
$\phi\ge\phi_{\mathrm{crit}}$. El $v_{\mathrm{in}}$ que cierra el
problema se busca por bisección (y, más adelante, una corrección tipo
secante). El operador directo es
\begin{equation}
\label{eq:F}
\mathcal{F}(\theta)
=q^\ast(\theta)
=v_{\mathrm{in}}^\ast(\theta)\,\rho_{t,i}(\theta),
\end{equation}
donde $\theta$ reúne $R$, $\Delta P$, $c_0$, $T$, $\xi_0$ y, si se
abre, $R(\cdot)$.

\section{El DAE, escrito de una vez}

Juntando todo, el modelo estacionario es
\begin{equation}
\label{eq:sistema}
\begin{aligned}
\frac{\dd P}{\dd z}
&= \mathcal{P}(P,\phi,N,\xi;q),\\
\frac{\dd\phi}{\dd z}
&= \Phi(P,\phi,N,\xi;q),\\
\frac{\dd N}{\dd z}
&= \frac{1}{u_m}\Gamma_N(P,\phi,N,\xi),\\
\frac{\dd\xi}{\dd z}
&= \frac{1}{u_m}\Gamma_\xi(P,\xi),\\
0
&= u_m-\frac{q\bigl(1-n(P,\xi)\bigr)}{(1-\phi)\rho_m},\\
0
&= u_g-\frac{q\,n(P,\xi)}{\phi\,\rho_g(P)},
\end{aligned}
\end{equation}
con $\mathcal{P}$ y $\Phi$ dados por~\eqref{eq:22}, $\Gamma_N$
por~\eqref{eq:dNdt} y $\Gamma_\xi$ por~\eqref{eq:dxdt}. Las fuerzas y
la viscosidad que entran en $\mathcal{P}$ y $\Phi$ dependen del
régimen $\phi\lessgtr\phi_1,\phi_2,\phi_{\mathrm{crit}}$. El dato de
borde~\eqref{eq:bc-base}--\eqref{eq:bc-top} determina $q$.

Eso es el modelo. El recorte monofásico de
\texttt{tesis/prototipo/} se obtiene de aquí al imponer
$u_m=u_g$, $\Gamma_N=\Gamma_\xi=0$,
$n=\max\bigl(0,\,(c_0-C_1P^\beta)/(1-C_1P^\beta)\bigr)$
y cortar el dominio en la fragmentación. No es otro volcán.

\section{Qué se recupera: el problema inverso}

Los observables de superficie (tasa de descarga, y si se quiere
$\phi(0)$, $u_g(0)$, $z_f$) son funcionales de la solución
de~\eqref{eq:sistema}. Invertir es, dado $d^{obs}$, hallar $\theta$
tal que $\mathcal{F}(\theta)\approx d^{obs}$. Las tres preguntas son
unicidad, estabilidad y reconstrucción. En el recorte
$u_m=u_g$ ya se demuestra que $\mathcal{F}$ no es inyectivo en
$R(\cdot)$. La pregunta de la tesis es si esa degeneración
\textbf{sobrevive} a $u_g\neq u_m$, que es precisamente lo que
~\eqref{eq:mom-m}--\eqref{eq:mom-g} tienen y el recorte no.

\section{La versión transiente, en breve}

\texttt{RIconduit1D\_transient.py} transporta las variables internas
por el conducto,
\begin{equation}
\label{eq:trans}
\frac{\p\phi}{\p t}+u_m\frac{\p\phi}{\p z}=S_\phi,
\qquad
\frac{\p N}{\p t}+u_m\frac{\p N}{\p z}=\Gamma_N,
\qquad
\frac{\p\xi}{\p t}+u_m\frac{\p\xi}{\p z}=\Gamma_\xi,
\end{equation}
y deja la presión en equilibrio hidrostático--viscoso de la mezcla,
$\dd P/\dd z=-\rho_{\mathrm{mix}}g-c_g\mu u_{\mathrm{mix}}/R^2$. Antes
de fragmentar usa una sola velocidad (HEM); después, un drift-flux de
Zuber--Findlay con $v_{\mathrm{slip}}$ de Newton. No es el mismo
cierre de momentum que~\eqref{eq:mom-m}--\eqref{eq:mom-g}: es una
variante, no un reemplazo. Un operador inverso serio tiene que
elegir \emph{uno} de los dos y quedarse con él.

\section{Constantes del caso Calbuco}

\begin{center}
\small
\begin{tabular}{@{}llr@{}}
\toprule
símbolo & significado & valor \\
\midrule
$H$ & techo de cámara & $-7000\,\mathrm{m}$ \\
$R$ & radio & $16\,\mathrm{m}$ \\
$\Delta P$ & sobrepresión de entrada & $5\times10^6\,\mathrm{Pa}$ \\
$c_0$ & agua total & $0.04$ \\
$T$ & temperatura & $1243.15\,\mathrm{K}$ \\
$\xi_0$, $\xi_{\max}$ & cristales inicial / techo & $0.25$, $0.50$ \\
$\tau_c$ & tiempo de cristalización & $2\,\mathrm{h}$ \\
$C_1$, $\beta$ & Henry & $4.11\times10^{-6}$, $1/2$ \\
$\Rv$ & constante del vapor & $461.11\,\mathrm{J\,kg^{-1}K^{-1}}$ \\
$\rho_{\mathrm{crust}}$ & densidad cortical & $2600\,\mathrm{kg\,m^{-3}}$ \\
$g$ & gravedad & $9.81\,\mathrm{m\,s^{-2}}$ \\
\bottomrule
\end{tabular}
\end{center}

\begin{thebibliography}{9}\small
\bibitem{bottinga1970} Bottinga, Y.\ \& Weill, D.\,F. (1970).
  \emph{Densities of liquid silicate systems calculated from partial
  molar volumes of oxide components}. Am.\ J.\ Sci.\ 269:169--182.
\bibitem{giordano2008} Giordano, D., Russell, J.\,K.\ \& Dingwell,
  D.\,B. (2008). \emph{Viscosity of magmatic liquids: a model}.
  EPSL 271:123--134.
\bibitem{kozono2009} Kozono, T.\ \& Koyaguchi, T. (2009).
  \emph{Effects of relative motion between gas and liquid on
  1-dimensional steady flow in silicic volcanic conduits}.
  JVGR 180:37--48.
  \href{https://doi.org/10.1016/j.jvolgeores.2008.11.006}{10.1016/j.jvolgeores.2008.11.006}
\bibitem{llewellin2006} Llewellin, E.\,W.\ \& Manga, M. (2005).
  \emph{Bubble suspension rheology and implications for conduit flow}.
  JVGR 143:205--217.
\bibitem{slezin2003} Slezin, Yu.\,B. (2003).
  \emph{The mechanism of volcanic eruptions (a steady state approach)}.
  JVGR 122:7--50.
\bibitem{yoshida1999} Yoshida, S.\ \& Koyaguchi, T. (1999).
  \emph{A new regime of volcanic eruption due to the relative motion
  of magma and liquid}. JVGR 89:215--226.
\end{thebibliography}

\end{document}
